\documentclass[a4paper]{article}
\usepackage{MMVII}

\newcommand{\TheCmd}{ImageMetadata}

\begin{document}

\CmdTitle
%\maketitle

\tableofcontents


 %=================================================================================

\SecDescr

This command print on the console  the meta-data 
on images that \MMVII can extract.  
Meta-data on images are used at different step of images calculation :

\begin{itemize}
   \item as initial value for computing the focal length in pixel;
   \item to decide which images correspond to the same camera,
   \item in radiometric processing
\end{itemize}


These meta data can come from :

\begin{itemize}
   \item value stored in the  \FileSty{xif} part of the  file
   \item rules added by user with command  \FileSty{EditCalcMTDI};
\end{itemize}

 %=================================================================================

\section{Display of values}

The printing depend of the different origin of the value :

\begin{itemize}
      \item case value \emph{exist} in xif and is \emph{not specified}  by \MMVII : print the value ;
      \item case value \emph{does not exist} in xif and is \emph{not specified} by \MMVII : print \FileSty{null};
      \item value  \emph{does not exist}   in xif and  \emph{is specified} bt \MMVII : print the value
            and indicate \FileSty{MMVII, SetByUserRule}
      \item value  \emph{exist}   in xif and  \emph{is modified} by \MMVII : 
             indicate with \emph{warning} color \FileSty{MMVII, ModifiedByUserRule},
            print the initial and modified values;
\end{itemize}


The otionnal parameter  \FileSty{Disp} control which meta-data information 
are printed, its default value  \FileSty{MainExif} prints the tags possibly involved in photogrammetric
or phototometric processing. 
The value  \FileSty{PhgrExif} prints only the tags used in photogrammetric computation.

If there is multiple file, the command groups the images by sets having exactly the same values
in all the meta-data specfied by  \FileSty{Disp}. For example with the ramses data set:

\begin{itemize}
   \item  \FileSty{MMVII ImageMetadata .*JPG } will generate $9$ group because there is different values of
          exposure time;
   \item  \FileSty{MMVII ImageMetadata .*JPG  Disp=PhgrExif} will generate $1$ group because all the images
          are acquired with the same camera;
\end{itemize}

 %=================================================================================

\section{Values specific to \MMVII}

There is values,  specific to \MMVII, usefull in processing that are not part
of exif information. When they exist they are printed :

\begin{itemize}
    \item if the  camera is in the data base (see  \FileSty{ EditCamDataBase}), the size
            of the sensor, and the size of the pixel are printed;

    \item if the  camera identifier is known it is printed, this identifier depend
          of the focal length, the camera model and the possible additionnal name created
          with \FileSty{AdditionalName} in   \FileSty{EditCalcMTDI};
\end{itemize}

Note that when merging the images having the same meta data, when it 
exist, the camera identifier is used to decided is two images belong to the same subset.
For example with polygon data set :

\begin{itemize}
    \item \FileSty{  MMVII ImageMetadata C.*JPG Disp=PhgrExif }
    \item all images have same focal and same camera model, but there is two
          camera, they have been  distinguished  using \FileSty{EditCalcMTDI},
          and we have two different identifier of camera  

    \begin{itemize}
             \item \FileSty{CalibIntr\_CamILCE-6400\_AddC1\_Foc16000}
             \item \FileSty{CalibIntr\_CamILCE-6400\_AddC2\_Foc16000}
    \end{itemize}

     \item so we will print in $2$ groups;
\end{itemize}




 %=================================================================================


\SecRelatedCmd

\begin{itemize}
   \item  \FileSty{ EditCalcMTDI}
   \item  \FileSty{ EditCamDataBase}
\end{itemize}

 %=================================================================================

\end{document}
